\chapter{How to Use This Book}
\label{ch:00}

This book accompanies Stanford's CS336: Language Modeling from Scratch. It follows the course's arc---from tokenization to alignment---but provides deeper mathematical treatment, more worked examples, and self-contained prerequisite coverage.

\section{The ``From Scratch'' Philosophy}

When we say ``from scratch,'' we mean it. By the end of this book, you will have seen how to:
\begin{itemize}
  \item Define the language modeling problem using information theory
  \item Build a tokenizer from raw text
  \item Construct a Transformer architecture layer by layer
  \item Train it efficiently on GPUs with custom kernels
  \item Scale it using empirical scaling laws
  \item Curate and filter training data
  \item Align the model to follow instructions and reason
\end{itemize}

No black boxes. No ``just import this library.'' Every component is derived, explained, and connected to the underlying mathematics.

\section{How to Read This Book}

\subsection{Sequential Reading}
If you are new to language models, read Part~0 first to solidify prerequisites, then proceed through Parts~I--VI in order. This is the recommended path.

\subsection{``I Know ML'' Path}
If you are comfortable with probability, linear algebra, and PyTorch, skip Part~0 and start with Chapter~6 (The Language Modeling Problem). Refer back to Part~0 when needed.

\subsection{Systems Focus}
For a systems-oriented reading, focus on Chapters~14--18 after reviewing the Transformer basics in Chapters~9--11.

\subsection{Assignment Companion}
Each CS336 assignment maps to specific chapters:
\begin{center}
\begin{tabular}{lll}
\toprule
\textbf{Assignment} & \textbf{Topic} & \textbf{Chapters} \\
\midrule
A1 & Basics & 6--12 \\
A2 & Systems & 14--17 \\
A3 & Scaling & 19--21 \\
A4 & Data & 22--23 \\
A5 & Alignment & 24--26 \\
\bottomrule
\end{tabular}
\end{center}

\section{Notation Conventions}

We use the following conventions throughout the book:
\begin{itemize}
  \item \textbf{Scalars} are lowercase italic: $x, y, \alpha$
  \item \textbf{Vectors} are lowercase bold: $\vx, \vy$
  \item \textbf{Matrices} are uppercase bold: $\vW, \vQ, \vK, \vV$
  \item \textbf{Sets} are calligraphic: $\vocab$ (vocabulary), $\mathcal{D}$ (dataset)
  \item $\log$ always means natural logarithm unless otherwise noted
  \item Complete notation reference: see the Notation Table on page~\pageref{ch:00:sec:notation_table}
\end{itemize}

\section{Neuroscience Analogies}

Throughout this book, you will encounter purple-shaded \textbf{Neuroscience Analogy} boxes. These connect machine learning concepts to brain science to build intuition. Every analogy box has three parts:
\begin{enumerate}
  \item \textbf{Where the analogy helps}: the genuine structural or functional similarity
  \item \textbf{Where the analogy breaks}: the important differences
  \item \textbf{What intuition to keep}: the one-sentence takeaway
\end{enumerate}

These analogies are thinking tools, not scientific claims about how LLMs work.

\section{Notation Table}
\label{ch:00:sec:notation_table}

\begin{longtable}{lp{10cm}}
\toprule
\textbf{Symbol} & \textbf{Meaning} \\
\midrule
$T$ & Sequence length (context window) \\
$V$ or $|\vocab|$ & Vocabulary size \\
$d$ or $d_{\text{model}}$ & Model / embedding dimension \\
$d_k, d_v$ & Key and value dimensions per attention head \\
$d_{\text{ff}}$ & Feed-forward hidden dimension \\
$n_h$ & Number of attention heads \\
$n_{\text{kv}}$ & Number of key-value heads (for GQA) \\
$L$ & Number of Transformer layers \\
$N$ & Number of model parameters \\
$D$ & Dataset size (in tokens) \\
$B$ & Batch size \\
$\eta$ & Learning rate \\
$\vtheta$ & Parameter vector \\
$p(x)$ & True / empirical distribution \\
$p_{\vtheta}(x)$ or $q(x)$ & Model distribution \\
$H(p)$ & Entropy \\
$H(p, q)$ & Cross-entropy \\
$\KL{p}{q}$ & KL divergence \\
$\PPL$ & Perplexity \\
$\text{softmax}(\cdot)$ & Softmax function (row-wise for matrices) \\
$\vQ, \vK, \vV$ & Query, Key, Value matrices \\
$C$ & Compute budget (FLOPs) \\
$\MFU$ & Model FLOPs Utilization \\
\bottomrule
\end{longtable}
